\subsection{第二轮修订回答了什么，仍未回答什么}

第二轮新主测试的第一轮冻结方法均值为4.064米，修订方法为3.920米，差值0.145米的配对95\%区间为$[-0.043,0.334]$。该区间跨零，不能据此宣布常规混合总体上有清晰的改进证据。修订方法相对BIC的差值区间同样跨零。与静态WLS的40.99\%均值改善回答的是另一个基线问题，不能替代对更强方法的比较。

稀疏与压力分布提供更具体的结果：在$n=8$的运动边界格点中，误差从10.788米降至6.217米，改善4.572米的配对区间为$[2.635,6.654]$；完整压力总体改善1.161米，区间$[0.661,1.666]$；完整边界总体改善0.425米，区间$[0.156,0.705]$。这些是预登记总体或格点上的探索性区间，没有校正全部比较的多重性。$n=8$上恒速稳健基线的误差仍更低，且该观测数下静态场景反而由1.888米退化至2.746米，说明解除计数排除没有解决稳健运动模型选择的全部困难。$n=6,7$也仍受训练数量限制。

静态目标没有被修订消除。在新主测试的180个静态正常案例中，修订误差1.419米，WLS为1.093米，均值比为1.299，95\%区间$[1.165,1.446]$。不仅点估计超过1.10容忍线，区间也全部高于该线。边界静态总体的比值为1.212，区间$[1.117,1.327]$。因此不能将总体平均收益解释为“标准静态场景不劣化”；该要求仍是明确失败项。

成本结果也有条件。正式新主测试中，修订方法的本机平均耗时约16.92毫秒，第一轮冻结方法约20.14毫秒，BIC约11.19毫秒。减少完整拟合起点并没有删除新增的交叉验证工作；不同模块的工作量和耗时不严格等价，而且这些是单轮、非随机化执行顺序的机器测量，不是稳定加速比或等预算算法优越性证明。

\subsection{分组校准改善的是哪一种平衡}
在相同的800例校准和1200例独立覆盖案例上，90\%全局区域实际覆盖88.33\%，平均半径7.690米；分组区域覆盖88.17\%，平均半径7.247米。分组没有提高这次总体覆盖。它将实际输出静态组与恒速组的覆盖由84.27\%和94.96\%，调整到88.04\%和88.38\%，体现的是两个可观测组之间的再平衡。

95\%名义水平下，全局覆盖94.58\%、平均半径11.779米，分组覆盖95.00\%、平均半径10.802米。这个单次结果支持继续研究分组规则的效率，但不能推广为所有水平都更好，也不应在看过结果后只报告95\%水平。两种水平与两种规则均预先声明并完整保留。

真实运动异常组在90\%分组区域下的覆盖仅66.78\%，全局规则为64.80\%。这是因为“输出模型类别”与“真实污染运动场景”不同，混合场景内的异质性没有被两个可观测组完全解释。拒绝后的子集与偏移分布结果也只能按各自指标描述，不能从总体共形定理自动获得条件保证。

本次90\%分组区域的覆盖率Wilson95\%区间为$[86.22\%,89.87\%]$，其上端低于90\%。这提示该固定校准规则的覆盖可能偏低，但不构成对共形秩定理的反例：定理的概率还对随机校准分数取平均，而本区间只描述固定校准后独立测试抽样的不确定性。报告同时保留了组样本量、秩、阈值与逐例分数，以便检查这种差异。

\subsection{过程与计算结果的独立核对}
两轮事件、方法注册、原始模型响应、工具证据和冻结源码均通过完整历史链审计，未报告不一致。第二轮独立统计脚本从保存的输入真值与输出位置重新计算距离、分场景统计、配对区间、有限样本校准秩、覆盖和拒绝结果，不调用定位器、不生成新数据，也不重新选择方法。123279项一致性检查通过，覆盖15860次实际预测、155组配对bootstrap比较及12个覆盖总体。原始核查输出保存在\path{verification/revision_results_check.json}；它检验记录与数值一致性，不能替代对数学假设或实验外推的判断。

第一轮冻结后计算为6192次，第二轮为15860次，合计22052次。第二轮声明的10个方法名称对应9个不同有效配置，四折消融是所选配置的准确别名，复用预测且明确标记。别名没有被计作额外独立计算；这也避免用相同预测的不同名称夸大实验工作量。另提供的新观测接口演示不计入以上正式评测次数。
